\section{Step 50 Pass the inverse accuracy gate}\label{h263-transform-s10}
After finite checks pass, run the complete instructor accuracy gate on the target
inverse. Retain its data hashes, rounding convention, byte order, mode and per-cohort
metrics. Compare all five limits and exact zeros. An acceptable peak error does
not establish an acceptable mean error or bias.

Measure transform work separately from the host accuracy runner. The runner
includes reference data preparation, script compilation, state edits and many
function calls. For entry-to-return instruction work, include validation, staging,
both passes, reduction, clipping and publication. Do not count only the vector
dot products and call that the cost of a whole transform.

\paragraph{Observe} All eight body/mode/byte-order combinations passed the six
prescribed cohorts. Different valid output rounding choices still need their own
accuracy check before entering the reconstruction loop.
\Needspace{8\baselineskip}
\paragraph{Check} Which claim follows from passing this gate?
\begin{itemize}
  \item[A.] The complete H.263 encoder is already interoperable.
  \item[B.] The tested inverse meets the measured Annex A accuracy requirements.
  \item[C.] Quantization is lossless for every video.
\end{itemize}
